\section{Conclusão}

As entalpias molares obtidas por calorimetria foram $\Delta H_A = \SI{-33,4}{kJ.mol^{-1}}$ para a dissolução do $\mathrm{NaOH}$ sólido, $\Delta H_B = \SI{-86,3}{kJ.mol^{-1}}$ para a reação do $\mathrm{NaOH}$ sólido com $\mathrm{HCl}$ aquoso e $\Delta H_C = \SI{-49,9}{kJ.mol^{-1}}$ para a neutralização entre as soluções de $\mathrm{HCl}$ e $\mathrm{NaOH}$. As três reações mostraram-se exotérmicas, como esperado, e os valores ficaram abaixo dos de referência em módulo, com desvios de \SI{25}{\%}, \SI{14}{\%} e \SI{11}{\%}, respectivamente, o que é coerente com a absorção de calor pelo calorímetro e com as perdas para o ambiente, desprezadas no cálculo.

Quanto à Lei de Hess, a soma $\Delta H_A + \Delta H_C = \SI{-83,3}{kJ.mol^{-1}}$ diferiu de $\Delta H_B$ em apenas \SI{3,5}{\%}. Como os erros sistemáticos foram semelhantes nas três determinações, eles praticamente se compensaram na comparação, e a relação $\Delta H_B = \Delta H_A + \Delta H_C$ foi verificada experimentalmente.

As principais fontes de erro foram o calor absorvido pelo calorímetro e perdido para o ambiente, a ausência da série de leituras de temperatura e da extrapolação na etapa A, o caráter higroscópico do $\mathrm{NaOH}$ sólido, a retenção de parte da solução de $\mathrm{NaOH}$ na proveta na etapa C e a resolução do termômetro. Apesar dessas limitações, o método permitiu determinar as entalpias com a ordem de grandeza correta e verificar a Lei de Hess, cumprindo os objetivos propostos.
